\section{Introduction}
There are many high-level data-analysis languages such as R, Python, MATLAB, and Julia which offer a sophisticated range of features for statistical tests and visualization~\cite{hadley2016r,embarak2018data,DBLP:conf/sc/PrabhuJRZHKJLGB11}. Additionally, these languages benefit from notebook environments such as Jupyter Notebook~\cite{DBLP:journals/cse/GrangerP21} or Quarto~\cite{Allaire_Quarto_2024}, which support an interactive exploration and manipulation of data. These technologies support domain experts in conducting their analyses without training in software engineering principles~\cite{DBLP:journals/rjour/Vidoni21}. 
In practice, however, analysis scripts~\cite{trisovic_largescale_2022,DBLP:conf/msr/SihlerPSTDD24}
and notebooks~\cite{DBLP:conf/msr/IslamAW24,DBLP:conf/msr/PimentelMBF19,10.1093/gigascience/giad113} suffer issues in reproducibility and structure~\cite{DBLP:conf/vl/0002HB20}. 
Various studies show that they lack executability~\cite{trisovic_largescale_2022,DBLP:conf/msr/PimentelMBF19,10.1093/gigascience/giad113,DBLP:conf/msr/IslamAW24} and hence reproducibility. These studies reveal that~\qtyrange{65}{77}{\percent} of the analyzed artifacts fail to execute due to missing dependencies, wrong execution order, or bad coding practices~\cite{trisovic_largescale_2022}. Moreover, research on code reuse in data science finds that~\qtyrange{70}{80}{\percent} of the analyzed code snippets are clones of other snippets with~\qty{50}{\percent} containing no unique code at all~\cite{10.7717/peerj.13933,DBLP:journals/programming/KallenW21}. Code reuse is even more prevalent across different projects. This indicates significant reproducibility issues of data science code.
Such artifacts are used frequently to help other researchers understand the applied methods and the steps conducted in the accompanied publication or to reuse parts within their own related analyses~\cite{DBLP:conf/vl/KoenzenES20,10.7717/peerj.13933}.
Reasons for these difficulties are manifold. Ranging from finding the used versions of packages and the language, to external dependencies like resource files. Additionally, they may rely on previous steps or expect the dataset to have a specific shape~\cite{trisovic_largescale_2022}.
While there is already some research on improving the reproducibility of data science scripts and notebooks, there is not yet a unified view on the various challenges that arise when trying to reproduce, replicate, or reuse these artifacts.
In this work, we address this by framing these difficulties as implicit assumptions, similar to classical software development~\cite{mamun2011review}. Additionally, we propose the systematic use of static analysis techniques to identify and document these implicit assumptions in the form of constraints, e.g., by using abstract interpretation to infer the values of variables~\cite{DBLP:conf/popl/CousotC77,10.1145/3689609.3689996}.
Such an analysis may happen either in the context of the original analysis environment~(cf.~\cite{donat2025r4r}) or retroactively when a researcher tries to reproduce, replicate, or reuse the artifact.
The results can then be used to add checks to the program that serve as an explicit documentation and as a dynamic validation for each run of the script.
Within the rest of this paper, we use the R~programming language as a primary example to illustrate specific cases. Furthermore, we use it as a basis for an explorative proof-of-concept implementation of our idea on top of \textit{flowR}, a static program analysis framework for~R~\cite{floopsla}. However, the proposed approach and most of the challenges are inherent to data analysis and hence language-agnostic. Therefore, we include a brief overview of the generalizability of the approach. Finally, we conclude with the implications of our work as well as a list of next steps in \cref{sec:conclusion,sec:future-plans}.
